%%%PAGE 99 rot=0 legibility=good summary=六道带草图的小题：①某车由静启动(a-x图像求速度与时间)、②圆形区域电场W_AC、③弹性绳E-h图像、④洗衣机ω≥√(g/R)、⑤变长度棒1/v-x图像与功、⑥万有引力F-h图像
%%%BLOCK id=099-01 ch=01 sec=a-x图像·非匀变速 kind=problem alt=none cont=none y=0.03-0.28
\textbf{某车由静启动}

①
%FIG p099_a 0.082 0.105 0.268 0.207
\notefig[0.25]{figures/p099_a.png}

$V^2=2ax$\quad $\dfrac{V^2}{2}=ax=$\ $a$-$x$图面积为$\dfrac{V^2}{2}$
% 原稿“=”与“a-x图”之间有涂改痕迹
\[ \therefore\ \frac{V_1^2}{2}=\frac{a_0x_1}{2}\qquad x_1\text{时}\ V_1=\sqrt{a_0x_1}\]
\[ t_1=\frac{\sqrt{a_0x_1}}{a_0}=\sqrt{\frac{x_1}{a_0}} \]
% CHECK: t_1 手写按匀加速 V_1/a_0 计算，而 0~x_1 段加速度并非恒定
\[ \frac{V_2^2}{2}=\frac{(x_2+x_2-x_1)a_0}{2}\qquad x_2\text{时}\ \therefore V_2=\sqrt{(2x_2-x_1)a_0}\]
\[ t_2=\frac{V_2-V_1}{a_0}=\sqrt{\frac{2x_2-x_1}{a_0}}-\sqrt{\frac{x_1}{a_0}} \]
%%%END
%%%BLOCK id=099-02 ch=09 sec=圆形区域电场·动能最大点 kind=problem alt=none cont=none y=0.24-0.34
②
%FIG p099_b 0.09 0.256 0.25 0.338
\notefig[0.25]{figures/p099_b.png}

在C动能最大，在C处做切线，电场指向C\\
垂直于C处切线
\[ W_{ac}=2qER\cos^{2}\alpha \]
%%%END
%%%BLOCK id=099-03 ch=06 sec=弹性绳·E-h图像 kind=problem alt=03 cont=none y=0.34-0.42
③
%FIG p099_c 0.095 0.338 0.255 0.41
\notefig[0.25]{figures/p099_c.png}

弹性绳\quad $L_0$原长

%FIG p099_d 0.295 0.328 0.45 0.412
\notefig[0.25]{figures/p099_d.png}

\red{\struck{$F-mg=ma$}}
% 原稿红笔所写，被红线划去；首字母 F/T 看不清（疑为 F）
%%%END
%%%BLOCK id=099-04 ch=04 sec=竖直面匀速圆周运动·洗衣机 kind=problem alt=none cont=none y=0.42-0.54
④
%FIG p099_e 0.055 0.425 0.238 0.528
\notefig[0.25]{figures/p099_e.png}

要衣物\unclear{贴}内壁则，$\omega\ge\sqrt{\dfrac{g}{R}}$

匀速圆周运动

洗衣机
%%%END
%%%BLOCK id=099-05 ch=12 sec=电磁感应·变长度棒与1/v-x图像 kind=problem alt=01,06 cont=none y=0.55-0.80
⑤
%FIG p099_f 0.025 0.572 0.23 0.668
\notefig[0.28]{figures/p099_f.png}

%FIG p099_g 0.24 0.563 0.445 0.68
\notefig[0.28]{figures/p099_g.png}

由$\dfrac{1}{V}\propto x$\ 为\uline{非匀变速运动}

\red{\uwave{图像面积为运动时间}}\quad $t=\dfrac{3L_0}{2V_0}$
\[ \frac{1}{V}=\frac{1}{L_0V_0}\cdot x+\frac{1}{V_0}\ \to\ V=\frac{L_0V_0}{x+L_0} \]
\[ E=B(L_0+x)V\qquad I=\frac{E}{R}=\frac{BL_0V_0}{R}\ \text{（定值）} \]
\[ W=I^2Rt-\Delta E_k=\frac{3B^2L_0^2V_0}{2R}-\frac{3mV_0^2}{8} \]
% CHECK: 分子中 L_0 的指数手写为 2，按 I^2Rt 推算似应为 L_0^3
%%%END
%%%BLOCK id=099-06 ch=05 sec=万有引力·F-h图像 kind=problem alt=none cont=none y=0.84-0.95
⑥
\[ F=\frac{GmM}{(R+h)^2} \]
%FIG p099_h 0.186 0.85 0.35 0.94
\notefig[0.25]{figures/p099_h.png}
%%%END
%%%PAGEEND 99
